% from ega0-22-fr.tex
\oldpage[0\textsubscript{IV-1}]{202}

  \item[\rm(iii)] Supposons que $A$ et $A'$ vérifient les hypothèses générales
  de (22.4.1). On sait (19.8.10) qu'il existe un anneau local \emph{régulier}
  $B$, d'idéal maximal $\mathfrak n$, tel que $A$ soit isomorphe à
  $B/\mathfrak n^2$~; désignons par $G_i\in B[T_i]$ un polynôme unitaire dont
  l'image canonique dans $A[T_i]$ soit $F_i$ $(1\leqslant i\leqslant r)$~;
  alors, si $B'$ est l'anneau quotient de $B[T_1,\ldots,T_r]$ par l'idéal
  engendré par les $G_i$, il est clair que $B'$ est un $B$-module libre de rang
  $d$ et que $A'=B'\otimes_BA=B'/\mathfrak n^2B'$~; comme $A'$ est supposé être
  un anneau local, il en est de même de $B'$, et si $\mathfrak n'$ est l'idéal
  maximal de $B'$, $B'/\mathfrak n'$ est isomorphe à $K'$ et
  $\mathfrak n'/\mathfrak n'^2$ isomorphe à $\mathfrak m'/\mathfrak m'^2$ en
  tant qu'espace vectoriel sur $K'$~; l'inégalité (22.5.2.1) montre donc que
  l'on a, sous les hypothèses de (22.4.1),
  \begin{equation}
  \label{0.22.5.2.2-fr}
    \operatorname{rg}_K(V)\leqslant\operatorname{rg}_{K'}(\mathfrak m'/\mathfrak m'^2).
    \tag{22.5.2.2}
  \end{equation}
\end{enumerate}
\end{remark}

\begin{corollary}[22.5.3]
\label{0.22.5.3-fr}
Soient $A$ un anneau local régulier, $\mathfrak m$ son idéal maximal,
$T_i$ $(1\leqslant i\leqslant r)$ des indéterminées, et pour chaque $i$, soit
\begin{equation}
\label{0.22.5.3.1-fr}
  F_i(T_i)=T_i^{d_i}+\sum_{1\leqslant k\leqslant d_i}c_{ik}T_i^{d_i-k}
  \tag{22.5.3.1}
\end{equation}
un polynôme unitaire de $A[T_i]$, tel que $d_i\geqslant2$ pour tout $i$ et
$c_{ik}\in\mathfrak m$. Soit $A'$ le quotient de l'anneau de polynômes
$A[T_1,\ldots,T_r]$ par l'idéal engendré par les $r$ polynômes $F_i$. Alors~:
\begin{enumerate}
  \item[\rm(i)] $A'$ est un anneau local, et son corps résiduel $K'$ est
  isomorphe au corps résiduel $K$ de $A$.
  \item[\rm(ii)] Pour que $A'$ soit un anneau régulier, il faut et il suffit
  que les classes $\xi_i$ mod.~$\mathfrak m^2$ des éléments $c_{i,d_i}$
  $(1\leqslant i\leqslant r)$ soient linéairement indépendantes sur $K$.
\end{enumerate}
\end{corollary}

\begin{proof}
En effet, avec les notations de (22.5.2, (i)), si $\overline F_i$ désigne le
polynôme de $A_1[T_i]$ obtenu en appliquant aux coefficients de $F_i$
l'homomorphisme $A\longrightarrow A_1$, $A'_1$ se définit à partir de $A_1$ et
des polynômes $\overline F_i$ comme dans (22.4.1), et la conclusion résulte
donc de (22.5.1) et de (22.4.3).
\end{proof}

\begin{theorem}[22.5.4]
\label{0.22.5.4-fr}
Soient $A$ un anneau local régulier, $K$ son corps résiduel~; on suppose que
$K$ est de caractéristique $p>0$~; soient
\begin{equation}
\label{0.22.5.4.1-fr}
  F_i(T_i)=T_i^p+p\sum_{k=1}^{p-1}c_{ik}T_i^{p-k}-f_i
  \tag{22.5.4.1}
\end{equation}
avec $c_{ik}\in A$ et $f_i\in A$ $(1\leqslant i\leqslant r)$~; soit $B$
l'anneau quotient de $A[T_1,\ldots,T_r]$ par l'idéal engendré par les $r$
polynômes $F_i$. Alors $B$ est un anneau local, et les conditions suivantes
sont équivalentes~:
\begin{enumerate}
  \item[a)] $B$ est régulier~;
  \item[b)] les éléments $d_Af_i\otimes1_K$ $(1\leqslant i\leqslant r)$ sont
  linéairement indépendants dans le $K$-espace vectoriel
  $\Omega^1_A\otimes_AK$~;
  \item[c)] si $B'=B/\mathfrak m^2B$ et si l'on pose
  $\Xi_{B_{(K)}/K}=\Xi_{B'_{(K)}/K}$ (22.4.5), l'homomorphisme
  caractéristique $\Xi_{B_{(K)}/K}\longrightarrow\mathfrak m/\mathfrak m^2$
  (22.4.6.2) est injectif.
\end{enumerate}
\end{theorem}

\begin{proof}
En effet, $B'$ se définit encore à partir de $A_1=A/\mathfrak m^2$ et des
polynômes $\overline F_i$ obtenus en appliquant aux $F_i$ l'homomorphisme
$A\longrightarrow A_1$. Compte tenu de ce que $\Omega^1_A\otimes_AK$ est
canoniquement isomorphe à $\Omega^1_{A_1}\otimes_{A_1}K$ (20.5.12, (ii)), le
théorème résulte de (22.5.1) et de (22.4.7.4), en utilisant la remarque
(22.4.8) lorsque $A_1$ n'est pas de caractéristique $p$.
\end{proof}

